\chapter{Dernières lignes peu ou prou communes}
\begin{mintedbox}{python}
"""
SOME CODE
"""


# création figure + axes
figure_ref_terre, axes_simu = plt.subplots()
figure_data, (axes_vitesse, axes_braquage) = plt.subplots(2, 1, sharex=True)
# titre
figure_ref_terre.suptitle("Animation")
figure_data.suptitle("Données en temps réel")
axes_vitesse.set_title("Vitesse")
axes_braquage.set_title("Braquage")
axes_braquage.set_xlabel("Temps")

# création de la forme de la voiture
shape_voiture = matplotlib.patches.Rectangle(POSITION_VOITURE_INIT, LONGUEUR_VEHIC, LARGEUR_VEHIC, fill=False, zorder=float('inf'))
shape_voiture.rotation_point = 'center'
shape_voiture.set_angle(INCLINAISON_VOITURE_INIT)

# création de la voiture
voiture = Voiture(shape_voiture, [axes_simu, axes_vitesse, axes_braquage], Tentacule_clothoïdaire)

# ajout du dessin de la voiture à la figure
axes_simu.add_patch(voiture.shape_voiture)

# création de la forme de l'obstacle
shape_obstacle = matplotlib.patches.Rectangle((30, -2.5), 2, 1.5, fill=True, zorder=float('inf'), color='cyan')
obstacle1 = Obstacle_fixe(shape_obstacle)
axes_simu.add_patch(obstacle1.patch)

# création des bords de route
shape_route_haut = matplotlib.patches.Rectangle((-10, 5), 90, 2)
route_haut = Obstacle_fixe(shape_route_haut)
axes_simu.add_patch(route_haut.patch)

shape_route_bas = matplotlib.patches.Rectangle((-10, -5), 78, 2)
route_bas = Obstacle_fixe(shape_route_bas)
axes_simu.add_patch(route_bas.patch)

shape_route_gauche = matplotlib.patches.Rectangle((68, -45), 2, 42)
route_gauche = Obstacle_fixe(shape_route_gauche)
axes_simu.add_patch(route_gauche.patch)

shape_route_droite = matplotlib.patches.Rectangle((78, -45), 2, 50)
route_droite = Obstacle_fixe(shape_route_droite)
axes_simu.add_patch(route_droite.patch)

# création de la simulation
simulation = Simulation_Ref_Terre([voiture], [obstacle1, route_bas, route_haut, route_droite, route_gauche])

# création et lancement de l'animation
anim = matplotlib.animation.FuncAnimation(figure_ref_terre,
                                          simulation.animer,
                                          init_func=lambda: [voiture.shape_voiture for voiture in simulation.voitures] + [tentacule.patch for tentacule in simulation.new_tentacules] + [tentacule.patch for tentacule in simulation.old_tentacules],
                                          fargs=[],
                                          interval=35,
                                          cache_frame_data=False,
                                          blit=True,
                                          repeat=True)

# création d'un timer
timer = figure_ref_terre.canvas.new_timer(interval=TEMPS_REAC)
if SECURITE is not False:
    timer2 = figure_ref_terre.canvas.new_timer(40)
    timer2.add_callback(simulation.collisions)
    timer2.start()
    simulation.ajout_timer(timer2)
    figure_ref_terre.canvas.mpl_connect('close_event', lambda event: timer2.stop())
# ajout d'une fonction appelée itérativement
timer.add_callback(simulation.gestion_tentac, voiture)
# démarrage du timer
timer.start()

simulation.ajout_anim(anim)
simulation.ajout_timer(timer)
# stop le timer si on ferme la figure (pour éviter qu'une fonction continue de tourner en boucle sans intérêt une fois tout terminé, c'est plutôt mal vu...!)
figure_ref_terre.canvas.mpl_connect('close_event', lambda event: timer.stop())
figure_ref_terre.canvas.mpl_connect('key_press_event', lambda event: (print("Key Pressed"), simulation.pause_replay()))

axes_simu.autoscale()
# axes_simu.set(xlim=(-50, 50), ylim=(-50, 150))
axes_simu.axis('equal')
plt.show()

\end{mintedbox}
